
\documentclass{article}
\usepackage{colortbl}
\usepackage{makecell}
\usepackage{multirow}
\usepackage{supertabular}

\begin{document}

\newcounter{utterance}

\centering \large Interaction Transcript for game `cladder', experiment `full\_v1.5\_default', episode 949 with llama{-}3.1{-}8b{-}instruct.
\vspace{24pt}

{ \footnotesize  \setcounter{utterance}{1}
\setlength{\tabcolsep}{0pt}
\begin{supertabular}{c@{$\;$}|p{.15\linewidth}@{}p{.15\linewidth}p{.15\linewidth}p{.15\linewidth}p{.15\linewidth}p{.15\linewidth}}
    \# & \multicolumn{2}{c}{Player} && \multicolumn{2}{c}{Game Master} \\
    \hline

    \theutterance \stepcounter{utterance}  
    & & & \multicolumn{4}{p{0.6\linewidth}}{
        \cellcolor[rgb]{0.9,0.9,0.9}{
            \makecell[{{p{\linewidth}}}]{
                \texttt{\tiny{[P1$\langle$GM]}}
                \texttt{You are an expert in causal inference. The following question is not a typical commonsense query, but rather a meticulously designed question created by a professor specializing in causal inference, intended to assess the students' mastery of the course content.} \\
\texttt{Imagine a self{-}contained, hypothetical world with only the following conditions, and without any unmentioned factors or causal relationships: Obesity has a direct effect on diabetes and lifespan. Smoking has a direct effect on diabetes and lifespan. Diabetes has a direct effect on lifespan. Smoking is unobserved. The overall probability of obesity is 10\%. For normal weight people, the probability of long lifespan is 86\%. For obese people, the probability of long lifespan is 54\%. Is long lifespan less likely than short lifespan overall?} \\
\texttt{Start your answer with 'Yes' or 'No', followed by additional reasoning or evidence to support your explanation.} \\
            }
        }
    }
     \\ \\

    \theutterance \stepcounter{utterance}  
    & \multicolumn{4}{p{0.6\linewidth}}{
        \cellcolor[rgb]{0.561,0.737,0.561}{
            \makecell[{{p{\linewidth}}}]{
                \texttt{\tiny{[P1$\rangle$GM]}}
                \texttt{No, long lifespan is not less likely than short lifespan overall.} \\
\\ 
\texttt{To answer this question, we need to apply the principles of causal inference, specifically the concept of counterfactuals and potential outcomes. Since smoking is unobserved, we cannot directly estimate its effect on lifespan. However, we can use the information provided to make an indirect inference.} \\
\\ 
\texttt{Let's define the following variables:} \\
\\ 
\texttt{{-} O: Obesity (observed)} \\
\texttt{{-} D: Diabetes (observed)} \\
\texttt{{-} L: Long lifespan (observed)} \\
\texttt{{-} S: Smoking (unobserved)} \\
\\ 
\texttt{We know that:} \\
\\ 
\texttt{{-} P(O) = 0.1 (probability of obesity)} \\
\texttt{{-} P(L|O) = 0.54 (probability of long lifespan given obesity)} \\
\texttt{{-} P(L|not O) = 0.86 (probability of long lifespan given normal weight)} \\
\\ 
\texttt{We also know that diabetes has a direct effect on lifespan, and obesity and smoking both have direct effects on diabetes and lifespan. However, since smoking is unobserved, we cannot directly estimate its effect on lifespan.} \\
\\ 
\texttt{To determine whether long lifespan is less likely than short lifespan overall, we need to consider the potential outcomes of being obese or normal weight, and the potential outcomes of having diabetes or not.} \\
\\ 
\texttt{Let's define the following potential outcomes:} \\
\\ 
\texttt{{-} P(L|O, D) = ? (probability of long lifespan given obesity and diabetes)} \\
\texttt{{-} P(L|not O, D) = ? (probability of long lifespan given normal weight and diabetes)} \\
\texttt{{-} P(L|O, not D) = ? (probability of long lifespan given obesity and no diabetes)} \\
\texttt{{-} P(L|not O, not D) = ? (probability of long lifespan given normal weight and no diabetes)} \\
\\ 
\texttt{Since diabetes has a direct effect on lifespan, we can assume that P(L|D) < P(L|not D). This means that having diabetes reduces the probability of long lifespan.} \\
\\ 
\texttt{Now, let's consider the potential outcomes of being obese or normal weight. We know that P(L|O) = 0.54 and P(L|not O) = 0.86. However, we also know that obesity increases the risk of diabetes, which reduces the probability of long lifespan.} \\
\\ 
\texttt{To determine whether long lifespan is less likely than short lifespan overall, we need to consider the potential outcomes of being obese or normal weight, and the potential outcomes of having diabetes or not. Since smoking is unobserved, we cannot directly estimate its effect on lifespan. However, we can use the information provided to make an indirect inference.} \\
\\ 
\texttt{Let's assume that the effect of obesity on diabetes is strong, and that the effect of diabetes on lifespan is also strong. In this case, the probability of long lifespan given obesity and diabetes (P(L|O, D)) would be lower than the probability of long lifespan given normal weight and no diabetes (P(L|not O, not D)).} \\
\\ 
\texttt{However, we also know that the probability of long lifespan given obesity (P(L|} \\
            }
        }
    }
    & & \\ \\

    \theutterance \stepcounter{utterance}  
    & & & \multicolumn{2}{p{0.3\linewidth}}{
        \cellcolor[rgb]{0.9,0.9,0.9}{
            \makecell[{{p{\linewidth}}}]{
                \texttt{\tiny{[GM$|$GM]}}
                \texttt{No} \\
            }
        }
    }
    & & \\ \\

    \theutterance \stepcounter{utterance}  
    & & & \multicolumn{2}{p{0.3\linewidth}}{
        \cellcolor[rgb]{0.9,0.9,0.9}{
            \makecell[{{p{\linewidth}}}]{
                \texttt{\tiny{[GM$|$GM]}}
                \texttt{no} \\
            }
        }
    }
    & & \\ \\

    \theutterance \stepcounter{utterance}  
    & & & \multicolumn{2}{p{0.3\linewidth}}{
        \cellcolor[rgb]{0.9,0.9,0.9}{
            \makecell[{{p{\linewidth}}}]{
                \texttt{\tiny{[GM$|$GM]}}
                \texttt{game\_result = WIN} \\
            }
        }
    }
    & & \\ \\

\end{supertabular}
}

\end{document}
